% GENERATED FILE. Edit canonical lessons, then run npm run generate:books.
% canonical-source-hash: fnv1a64:de2fcd2a29ef63a6
% canonical-lessons: ML-C25-shubha-rathri, ML-S129-letter-bha

\chapter{Good Night}
\label{ch:ml-good-night}
\hlchaptermodality{25}

\begin{chapteropening}
\textbf{By the end of this chapter:} \emph{I can wish someone good night in Malayalam with śubha rāthri.}

Say śubha rāthri as 'auspicious night', and mark the English code-switching claim as thinly sourced and not unique to Malayalam.
\end{chapteropening}

\section[śubha rāthri]{\ml{ശുഭ} \ml{രാത്രി} (śubha rāthri) — \textquotedblleft{}good night,\textquotedblright{} and a pattern seen before}

\label{lesson:ML-C25-shubha-rathri}



\begin{quote}
You already have \emph{rāthri}. One more Sanskrit word, shared with Hindi, Kannada, and Telugu, completes Malayalam's \textquotedblleft{}good night\textquotedblright{} — and this lesson's honest twist is one you've technically seen before, just in a different language.
\end{quote}

\subsection*{What to know first}
\textbf{\ml{ശുഭ} \ml{രാത്രി}} (\textbf{śubha rāthri}) — \textquotedblleft{}\textbf{good night}\textquotedblright{} — literally \textquotedblleft{}\textbf{auspicious night}.\textquotedblright{} Both pieces are Sanskrit \textbf{tatsama}: \textbf{\ml{രാത്രി}} (\emph{rāthri}, \textquotedblleft{}night,\textquotedblright{} already met last lesson) plus \textbf{\ml{ശുഭ}} (\emph{śubha}, \textquotedblleft{}\textbf{auspicious},\textquotedblright{} from root \textbf{śubh}, \textquotedblleft{}to be beautiful,\textquotedblright{} ultimately PIE \emph{\textbf{*\'{k}ewb\textsuperscript{h}-}} — the same word, the same root, already verified for the Hindi, Kannada, and Telugu siblings of this exact arc).

\begin{cousinweb}
Here's the honest twist, hedged deliberately: some sources describe modern Malayalam speakers, especially in casual settings, \textbf{code-switching to English} \textquotedblleft{}\textbf{good night}\textquotedblright{} rather than saying \textbf{\ml{ശുഭ} \ml{രാത്രി}} aloud — treating it more as the standard, correct-but-formal translation than an everyday spoken phrase. But this lesson isn't presenting that as a fresh Malayalam discovery. The \textbf{exact same claim}, sourced from similarly thin blogs and forum posts, is also made for \textbf{Telugu}, and there too it is genuinely \textbf{unclear} whether the pattern is Telugu-specific or something broader. Seeing the same claim turn up again for Malayalam \textbf{strengthens the suspicion that this is a broader South Asian pattern} (or simply a sourcing artifact — the languages that happen to get commented on in casual sources), not a language-specific trait worth treating as a genuine discovery each time. Either way, \emph{śubha rāthri} itself remains perfectly correct, understood Malayalam.
\end{cousinweb}

\subsection*{Your turn}
Say these aloud:

\begin{itemize}
\raggedright
  \item \textquotedblleft{}śubha rāthri\textquotedblright{} — \textquotedblleft{}good night,\textquotedblright{} literally \textquotedblleft{}auspicious night\textquotedblright{}
  \item śubha's own root — śubh, \textquotedblleft{}to be beautiful\textquotedblright{} $\to$ \textquotedblleft{}auspicious\textquotedblright{} — the same root already met for Hindi, Kannada, and Telugu
  \item the honest twist — this code-switching claim isn't uniquely Malayalam's; the same weakly-sourced claim already appeared for Telugu
\end{itemize}

\begin{tcolorbox}[breakable,colback=teal!4,colframe=teal!35!black,title={Before you move on}]
What does \textbf{\ml{ശുഭ} \ml{രാത്രി}} literally mean? (\textquotedblleft{}\textbf{Auspicious night}.\textquotedblright{}) Is the \textquotedblleft{}Malayalam speakers code-switch to English instead\textquotedblright{} claim a fresh discovery in this lesson? (\textbf{No} — the identical claim, similarly thinly sourced, is also made about \textbf{Telugu}, suggesting a broader pattern or a sourcing artifact, not a Malayalam-specific fact.) Does this affect whether \emph{śubha rāthri} itself is correct? (\textbf{No} — it's perfectly correct and understood Malayalam, whatever its actual spoken frequency.)
\end{tcolorbox}
\section[bha]{\ml{ഭ} — one character, met inside words you already say}

\label{lesson:ML-S129-letter-bha}



\begin{quote}
Before the new one: \ml{സ} — what does it do?

One character this time. One only. It has been sitting in front of you on pages you have already read.
\end{quote}

\subsection*{Script you'll notice: \ml{ഭ}}
\textbf{\ml{ഭ}} — \emph{bha}.

It is a \textbf{consonant}, and in this script a consonant is never bare: it comes with an \emph{a} already in it. So it is not \emph{bh}, it is \textbf{bha}.

Its sound is a \emph{b} with a puff of breath after it — \textbf{bha}, not \emph{ba}. Two words already in your hands are built on it, and both come from the Sanskrit layer, which is the pattern this letter keeps.

\begin{itemize}
\raggedright
  \item \textbf{\ml{ശുഭ}} \emph{śubha} — good, auspicious — the first half of the good-night wish on the page before this one
  \item \textbf{\ml{കുംഭം}} \emph{kuṁbhaṁ} — the month, met with the calendar
\end{itemize}

Say each one and listen for the breath: \emph{śu-bha}, \emph{kuṁ-bhaṁ}.

\subsection*{Writing: \ml{ഭ} — copy what you see}
Put your pen on \ml{ഭ} and follow its line. Copy the shape you can see — slowly, and larger than it is printed.

\begin{quote}
This book does not yet tell you \textbf{where to start the character or which way to travel}. That is a real thing, taught with real variation from school to school, and it is not written down here until it can be written down with a source. Copying what is in front of you needs no such source, and it is how the shape gets into your hand in the meantime.
\end{quote}

\subsection*{Your turn}
\begin{itemize}
\raggedright
  \item \emph{Look:} at these words, and find \ml{ഭ} in the ones that have it
\end{itemize}

\begin{quote}
\ml{ശുഭ}  ·  \ml{കുംഭം}  ·  \ml{ഭ}
\end{quote}

\begin{itemize}
\raggedright
  \item \emph{Trace it:} \ml{ഭ} three times, saying \emph{bha} as you finish each one
  \item \emph{Look:} back at any page of this chapter and find \ml{ഭ} once more
\end{itemize}

\begin{tcolorbox}[breakable,colback=teal!4,colframe=teal!35!black,title={Before you move on}]
Which character is this — \ml{ഭ}? What sound does it carry? (\emph{\textbf{bha}}.) Name one word you already say that contains it.
\end{tcolorbox}
